% Research-question-first body. Shared by the normal and 35-minute entry points.
\begin{frame}[plain]\titlepage\end{frame}

\begin{frame}{Start with a Task: Several Views, One Result}
\begin{columns}[T,onlytextwidth]
\column{.46\linewidth}
\begin{itemize}
\item An operator selects a ground area.
\item Nearby UAVs offer suitable views.
\item A compute Provider analyzes the images.
\item The operator receives original images with annotations.
\end{itemize}
\column{.51\linewidth}
\begin{tikzpicture}[node distance=8mm,box/.style={draw=MemphisBlue,rounded corners,fill=LightBand,align=center,font=\small,text width=26mm,minimum height=10mm},>=Stealth]
\node[box] (a) {UAV A\\view 1};
\node[box,below=of a] (b) {UAV B\\view 2};
\node[box,right=9mm of a,yshift=-9mm] (c) {Compute\\Provider};
\node[box,below=12mm of c] (u) {Operator};
\draw[->,thick] (a)--(c);\draw[->,thick] (b)--(c);\draw[->,thick] (c)--(u);
\end{tikzpicture}
\end{columns}
\vspace{3mm}\bluebox{Who can participate now, and which Data may each selected Provider use?}
\end{frame}

\begin{frame}{The Task Creates Three Requirements}
\begin{enumerate}
\item \textbf{Discover current willingness:} reachability, sensors, load, and deadlines affect participation.
\item \textbf{Authorize before execution:} exposing a service does not grant everyone permission to use it.
\item \textbf{Enforce the chosen collaboration:} an authentic image is useful only if it satisfies an assigned dependency.
\end{enumerate}
\vspace{4mm}\bluebox{The contribution must explain the mechanism and its cost, not just demonstrate that two UAVs can exchange images.}
\end{frame}

\begin{frame}{What NDN Provides---and What It Does Not}
\small
\begin{tabularx}{\linewidth}{p{.39\linewidth}X}
\toprule NDN building block & Remaining application responsibility\\\midrule
Named Data and signatures [1] & Decide which signer may perform which task; validate content under trust rules\\\midrule
Interest retrieval and caching [1,26] & Enforce task identity, dependencies, and deadlines; caching is not persistent storage\\\midrule
Dataset synchronization / SVS [2,3] & Discover new names, then retrieve Data; define service-message state transitions\\\midrule
NAC [4] / enhanced NAC-ABE [19] & Relate decryption permission to execution authorization and key lifecycle\\\bottomrule
\end{tabularx}
\vspace{2mm}\scriptsize NAC: Name-based Access Control; ABE: attribute-based encryption; SVS: State Vector Sync.
\vspace{2mm}\bluebox{\small NDNSF builds on these mechanisms; it does not invent them.}
\end{frame}

\begin{frame}{Three Research Questions}
\begin{itemize}
\item \textbf{RQ1:} How can a requester discover and authorize currently willing Providers without making the discovery descriptor public, and at what cost?
\item \textbf{RQ2:} How can Providers selected at runtime follow task roles and named dependencies, whether Data comes from producers or caches?
\item \textbf{RQ3:} When do these mechanisms improve completion, latency, or execution efficiency, and when do their costs outweigh the benefits?
\end{itemize}
\vspace{3mm}\bluebox{UAV and distributed inference test the same questions. They are not separate contributions merely because both are implemented.}
\end{frame}

\begin{frame}{Related Work Already Provides Important Pieces}
\small
\begin{tabularx}{\linewidth}{p{.27\linewidth}X}
\toprule Prior work & What the comparison must acknowledge\\\midrule
NFN [6], NFaaS [7] & Named computation and function execution/placement\\\midrule
RICE [9], NSC [10] & Invocation and security mechanisms; NSC validates signed calls using application trust rules\\\midrule
DNMP [12] & Role- and command-constrained signing and verification\\\midrule
CaSCON [15] & Service-composition research; requires a feature-level comparison before claiming a new composition mechanism\\\midrule
gRPC [17,27] & Complete configurations can include discovery/resolver integration, not only fixed addresses\\\bottomrule
\end{tabularx}
\vspace{2mm}\bluebox{Candidate increment: connect runtime assignment to enforceable object-admission rules. This is not a claim to be the first NDN service framework.}
\end{frame}

\begin{frame}{Contributions and Evidence Are Different Things}
\small
\begin{tabularx}{\linewidth}{p{.38\linewidth}X}
\toprule Intended contribution & Current evidence boundary\\\midrule
Discovery / authorization design and tradeoffs & Invocation implementation and bounded experiments exist; full alternative lifecycle-cost comparison remains open\\\midrule
ACK-derived plan and local dependency admission & Design and implementation work; complete direct/cache and failure acceptance gates remain open\\\midrule
Cross-case empirical findings & Existing mobility/work-efficiency and small-model results support limited configurations, not all new collaboration claims\\\bottomrule
\end{tabularx}
\vspace{3mm}\bluebox{Proposal status, September 2026: implemented, locally tested, measured, and fully qualified are not interchangeable. [28]}
\end{frame}

\begin{frame}{Participants and Trust Bootstrap}
\begin{itemize}
\item User: requests and coordinates the invocation.
\item Provider: offers and executes a service. A node may play both roles.
\item Controller: administers identities and service permissions approved by the organization.
\item Participants use provisioned trust anchors and validation rules for certificates and Controller material.
\end{itemize}
\vspace{3mm}\bluebox{Unknown runtime candidates are not unregistered authorization identities. The User discovers current participation; the Controller does not choose this task's Providers.}
\end{frame}

\begin{frame}{One Invocation, Four Named Data Types}
\small
\begin{tabularx}{\linewidth}{p{.15\linewidth}X}
\toprule Data & Logical purpose\\\midrule
Request & Describe required service and participation constraints; include protected discovery information and Requester Challenge\\\midrule
ACK & Report positive or negative participation; positive ACK supplies selection information and Provider Challenge\\\midrule
Selection & Authorize selected execution; provide input/references and protected request-scoped keys\\\midrule
Response & Return status, result Data, or authorized result references\\\bottomrule
\end{tabularx}
\vspace{3mm}\bluebox{Sync helps discover names; Interests retrieve the Data. A positive ACK does not itself start execution.}
\end{frame}

\begin{frame}{Descriptor Before Selection; Input After Selection}
\begin{columns}[T,onlytextwidth]
\column{.48\linewidth}
\begin{block}{Discovery descriptor}
Needed to decide whether to participate: area, deadline, sensor/model and resource requirements.
\end{block}
\column{.48\linewidth}
\begin{block}{Application input}
Needed to execute the selected task: images, detailed processing parameters, or protected input references.
\end{block}
\end{columns}
\vspace{4mm}\bluebox{The boundary follows necessity. A location needed to assess coverage belongs in the protected descriptor; an image needed only for processing can wait.}
\end{frame}

\begin{frame}{Reason 1: Confidential Discovery Before Selection}
\begin{itemize}
\item The User knows the service, but not yet the currently willing Provider identities.
\item An observation area and deadline can reveal intended activity even without the full application input.
\item ABE can protect one descriptor for a service-authorized receiver set without listing those identities. [19]
\end{itemize}
\vspace{3mm}\bluebox{Confidential from parties lacking the required decryption capability. Authorized but unselected Providers can read it; names and traffic patterns remain visible.}
\end{frame}

\begin{frame}{Reason 2: Separate Dissemination from Access}
\begin{itemize}
\item Nodes may retrieve the same encrypted Request Data; only matching decryption capabilities expose its descriptor.
\item The User encrypts by service attributes, without listing Provider identities. The Controller still authorizes entities and issues keys. [19]
\item Different entities retain overlapping service permissions in individual policy keys, independently of identity-signing credentials.
\end{itemize}
\vspace{3mm}\bluebox{Architectural benefit: shared dissemination with distinct access scopes. Aggregation preserves distinctions only while objects retain separate encryption scopes.}
\end{frame}

\begin{frame}{A Fair Comparison Includes Encryption on Both Sides}
\small
\begin{tabularx}{\linewidth}{p{.27\linewidth}XX}
\toprule Design & Execution authorization & Discovery confidentiality\\\midrule
DNMP-inspired roles [12] & Role/command signing rules, credential and request-state checks & Add encryption and group/key distribution\\\midrule
Our per-service role certificates & Service/message rules plus caller or Provider credential & Add equivalent encryption and lifecycle handling\\\midrule
NDNSF ABE-backed & Identity signature, encrypted challenges, and authorization/request state & Service-level ABE protection of descriptor Content\\\bottomrule
\end{tabularx}
\vspace{3mm}\bluebox{Trust Schema [16] validates signing relationships; it alone does not encrypt Content. NAC [4] also protects distributed content. ABE contributes attribute-policy-based decryption.}
\end{frame}

\begin{frame}{When Another Authorization Design May Be Better}
\begin{itemize}
\item Known Providers and simple rules: signed invocation plus recipient encryption may be sufficient.
\item A small, stable receiver group: ordinary group-key management may be attractive.
\item Rich command-argument restrictions: DNMP-style rules make those constraints explicit. [12]
\item Many changing entitlements: measure both policy/key updates and credential/rule updates before claiming a management advantage.
\end{itemize}
\vspace{3mm}\bluebox{ABE is a justified option for NDNSF's requirements, not universally superior authorization.}
\end{frame}

\begin{frame}{Authorization Lifecycle and Costs}
\begin{itemize}
\item Challenges check decryption capability; signatures, service permissions and request state still determine invocation acceptance.
\item Request-scoped keys protect input/results. Selection delivers keys to selected Providers; the requesting User retains access.
\item Withdrawal requires retrieved, validated Controller status. Current ABE rotation is Controller-wide. [28]
\end{itemize}
\vspace{3mm}\bluebox{Lifecycle requirement, not a third advantage: measure policy/key size, challenge cost, propagation and rekeying. No immediate disconnected revocation or recall of old plaintext/keys.}
\end{frame}

\begin{frame}{Threat Model: What Is Trusted?}
\begin{itemize}
\item Trusted: Controller authorization and the User's planning decision.
\item Untrusted delivery: an object may arrive through a producer, forwarder, cache, or repository and still requires validation.
\item Adversarial cases to test: wrong identities, permissions, challenges, request state, roles, dependencies, or execution identities.
\item Not established by signatures/ABE: honest resource reports, correct inference, truthful images, or resistance to authorized key sharing.
\end{itemize}
\bluebox{No Byzantine-coordinator consensus claim; no traffic-analysis or instantaneous-revocation guarantee.}
\end{frame}

\begin{frame}{RQ2: Generate the Plan from Eligible Positive ACKs}
\begin{enumerate}
\item Application supplies roles, dependency template, constraints and deadline; the template may be handwritten.
\item Close the ACK collection window and validate eligible positive ACKs.
\item A replaceable policy assigns Providers to roles and produces a signed plan.
\item Each Provider derives its permitted input edges and rechecks local resource admission before starting.
\end{enumerate}
\vspace{3mm}\bluebox{The snapshot freezes observed candidates for this decision, not every Provider in the network or their future availability.}
\end{frame}

\begin{frame}{A Plan Is More Than a Request ID}
\small
\begin{tabularx}{\linewidth}{p{.3\linewidth}X}
\toprule Existing request ID & Proposed task-specific plan rules\\\midrule
Binds messages to a request & Assigns exact Providers to collection and processing roles\\\midrule
Does not alone specify predecessors & Declares which producer role may supply each named dependency\\\midrule
Does not alone describe key scope & Associates permitted objects, execution identity, and authorization/key scope\\\bottomrule
\end{tabularx}
\vspace{4mm}\bluebox{Example: within the same request, a service-authorized UAV that was not assigned ``view 1'' must not satisfy the view-1 dependency.}
\end{frame}

\begin{frame}{Plan-to-Data Enforcement: Check Before Advancing}
\begin{enumerate}
\item Validate the Data signature using the appropriate trust rules.
\item Match the current request, plan and execution identity.
\item Verify that the signer is the assigned predecessor Provider.
\item Match the declared dependency edge, name, data schema and key scope.
\item Require complete validated input before the dependent computation starts.
\end{enumerate}
\vspace{3mm}\bluebox{Apply the same rules to direct, cached and repository-fetched Data. This is a proposed acceptance contract; full end-to-end qualification remains required.}
\end{frame}

\begin{frame}{Failure Handling Needs Explicit State}
\small
\begin{tabularx}{\linewidth}{p{.3\linewidth}X}
\toprule Event & Required behavior\\\midrule
Late ACK & Do not silently change the committed plan\\\midrule
Duplicate Selection & Follow recorded one-time state; do not implicitly repeat a side effect\\\midrule
Incomplete dependency & Wait within policy/deadline, recover, or fail; do not execute on partial data\\\midrule
Provider replacement & Authorize a new assignment; define fencing and whether existing results may be reused\\\midrule
Cancellation & Prevent late messages from reauthorizing cancelled work\\\bottomrule
\end{tabularx}
\vspace{3mm}\bluebox{Network re-retrieval is not task re-execution. Compensation is not an atomic transaction or exactly-once physical action.}
\end{frame}

\begin{frame}{Proposed UAV Flow: Capture, Analyze and Return}
\begin{enumerate}
\item Sensing UAVs publish encrypted ViewImage Data in their own authorized namespaces.
\item The compute Provider retrieves images over assigned edges and publishes Annotation Data.
\item The result references the original images and securely provides permitted image keys to the requester.
\item The requester retrieves the original images and displays the annotations.
\end{enumerate}
\vspace{3mm}\bluebox{No need to republish the same images as if the compute Provider produced them. GPS alone does not establish camera visibility or accurate triangulation.}
\end{frame}

\begin{frame}{What the UAV Case Must Demonstrate}
\begin{itemize}
\item A complete authorized collection--processing--return task before the deadline.
\item Correct handling of unavailable views, delayed images and a failed processing Provider.
\item Wrong-role/edge rejection without rejecting an authorized cached image.
\item MiniNDN [22] and PX4 SITL [23] results clearly distinguished from physical-flight evidence.
\end{itemize}
\vspace{3mm}\bluebox{Higher recognition/localization accuracy is a separate claim requiring ground truth and a single-view control. It is not assumed here.}
\end{frame}

\begin{frame}{Proposed DI Flow: Compute Follows Named Dependencies}
\begin{center}
\begin{tikzpicture}[box/.style={draw=MemphisBlue,fill=LightBand,rounded corners,text width=30mm,minimum height=12mm,align=center,font=\small},>=Stealth]
\node[box] (a) {Provider A\\layer group 1};
\node[box,right=17mm of a] (b) {Provider B\\layer group 2};
\node[box,right=17mm of b] (c) {Provider C\\output stage};
\draw[->,thick] (a)--node[above,font=\scriptsize]{activation} (b);
\draw[->,thick] (b)--node[above,font=\scriptsize]{activation} (c);
\end{tikzpicture}
\end{center}
\begin{itemize}
\item Pipeline: sequential layers/groups; tensor partitioning additionally needs supported collectives.
\item ONNX [25] describes the graph; export alone does not implement distributed execution.
\item Repo [24] stores artifacts; forwarding caches do not guarantee persistence.
\end{itemize}
\bluebox{The framework contribution is checked runtime assignment and data exchange, not model partitioning itself.}
\end{frame}

\begin{frame}{DI State and Evidence Boundaries}
\begin{itemize}
\item Model weights, reusable artifacts, request activations and KV state have different lifetimes.
\item Autoregressive decoding requires retained K/V and next-step inputs; resident weights do not remove that state dependency.
\item Native/local tests are not cross-process or multi-device qualification.
\item Freeze one supported model/partition first; tensor/hybrid and LLM streaming need their own evidence if claimed.
\end{itemize}
\vspace{3mm}\bluebox{Earlier small-model measurements do not qualify all current V3/native features. [28]}
\end{frame}

\begin{frame}{RQ3: Define the Baselines Before Comparing}
\begin{itemize}
\item NSC [10]: named service calls; describe the tested executor, timeout, queue and recovery behavior.
\item gRPC [17]: describe resolver, sequential/parallel policy, deadlines, retry and cancellation.
\item Match Provider count, capacity, task, topology/traces, publication times and measurement windows.
\item Report fixed-list and dynamically discovered configurations as different configurations.
\end{itemize}
\bluebox{An incomplete baseline cannot establish a limitation of the full framework. Historical unmatched data remain diagnostics.}
\end{frame}

\begin{frame}{Mobility: Overall Completion Depends on Coverage}
\small Ten seeds (62--71), four Providers/system, 2 m/s; mean seed completion rate. [28]
\vspace{4mm}
\begin{tabularx}{\linewidth}{Xrrr}
\toprule AP radius & NDNSF & gRPC & NSC\\\midrule
50 m & 54.57\% & 53.80\% & 55.47\%\\\midrule
100 m & 97.97\% & 97.97\% & 98.33\%\\\midrule
150 m & 100.00\% & 100.00\% & 100.00\%\\\bottomrule
\end{tabularx}
\vspace{4mm}
\bluebox{These overall controls do not establish NDNSF superiority. The next slide examines the separately registered switching subset, not all requests.}
\end{frame}

\begin{frame}{Mobility: Lower Tail Latency in a Defined Subset}
\small
Ten seeds; four Providers/system; 100 m, 2 m/s; 1,312 matched requests requiring a Provider switch. [28]
\vspace{4mm}
\begin{tabularx}{\linewidth}{Xrr}
\toprule Comparison & Seed-p95 difference & 95\% CI\\\midrule
NDNSF minus gRPC-SEQ-4 & $-597$ ms & $[-992,-200]$ ms\\\midrule
NDNSF minus NSC-4 & $-2618$ ms & $[-2916,-2315]$ ms\\\bottomrule
\end{tabularx}
\vspace{4mm}\bluebox{This is a conditional latency result, not general completion-rate superiority. Preserve the registered subset definition, overall coverage controls, and failed requests.}
\end{frame}

\begin{frame}{Selection Avoids Redundant Starts in This Comparison}
\small Six-seed, capacity-matched comparison; 1,800 requests per system. [28]
\vspace{4mm}
\begin{tabularx}{\linewidth}{Xrr}
\toprule Configuration & Completed requests & Execution starts\\\midrule
NDNSF & 1,798 / 1,800 & 1,800\\\midrule
gRPC-PAR-4 & 1,800 / 1,800 & 7,180\\\bottomrule
\end{tabularx}
\vspace{3mm}\small NSC-4 also completed 1,800/1,800, with 1.002 starts/request (rounded). No work-efficiency advantage over NSC is claimed.
\vspace{3mm}\bluebox{Fewer starts than gRPC-PAR-4 is not proof of fewer bytes, less CPU time or lower energy. Similar observed completion is not proven equivalence.}
\end{frame}

\begin{frame}{A Useful Negative Result: Splitting Can Be Slower}
\small Earlier small Qwen-style ONNX model; ten successful runs per layout. [28]
\vspace{4mm}
\begin{tabularx}{\linewidth}{Xrr}
\toprule Layout & Mean latency & p95\\\midrule
Shared-backbone distributed layout & 259.4 ms & 283.9 ms\\\midrule
Single Provider & 186.8 ms & 206.3 ms\\\bottomrule
\end{tabularx}
\vspace{4mm}\bluebox{Do not split merely because several Providers exist. Communication/setup overhead is a candidate explanation, not a causal result established by this measurement.}
\vspace{2mm}\scriptsize This campaign does not validate the newer canonical-layer, tensor/hybrid or multi-accelerator design.
\end{frame}

\begin{frame}{Remaining Experiments Follow the Claims}
\small
\begin{tabularx}{\linewidth}{p{.09\linewidth}XX}
\toprule RQ & Required evidence & Outcome and cost\\\midrule
RQ1 & Full authorization alternatives; unauthorized call, challenge/replay and withdrawal cases & No unauthorized handler start; discovery delay, key/policy size, status delay and rekey traffic\\\midrule
RQ2 & One UAV and one DI graph; wrong role/edge; direct/cache parity; cancellation/replacement & Correct task completion; validation, transfer and recovery cost\\\midrule
RQ3 & Matched workload and failure schedules; feasible single-device control; cold/warm state & Completion plus latency, redundant work, bytes, uncertainty, and cases with no advantage\\\bottomrule
\end{tabularx}
\vspace{2mm}\bluebox{Planned tests are not completed findings. Record all scheduled runs and failure boundaries.}
\end{frame}

\begin{frame}{A Bounded Path to Spring 2027}
\small
\begin{tabularx}{\linewidth}{p{.25\linewidth}X}
\toprule Internal target & Completion gate\\\midrule
September 2026 & Agree scope/RQs; freeze one UAV graph and one DI partition; resolve correctness blockers\\\midrule
October--November 2026 & Complete registered experiments; freeze evidence and write bounded RQ answers\\\midrule
December 2026 & Complete integrated advisor-review draft\\\midrule
January 2027 & Committee review and readiness decision; reduce scope or revise timing if core evidence is missing\\\midrule
February 2027 & Target defense, conditional on evidence and committee approval\\\bottomrule
\end{tabularx}
\vspace{2mm}\bluebox{CS guidance: advisor-approved complete draft to committee two weeks before defense. Spring 2027 final ProQuest deadline: March 26, 2027; recheck before filing. [29,30]}
\end{frame}

\begin{frame}{What a Final Defense Must Be Able to Claim}
\begin{itemize}
\item \textbf{RQ1 answer:} when ABE-backed discovery/authorization is preferable, including management and revocation costs.
\item \textbf{RQ2 answer:} exactly which task-assignment mistakes are rejected and which recovery behavior is supported.
\item \textbf{RQ3 answer:} measured benefits, negative results and deployment limits.
\item Distinguish the candidate's research from shared implementations and reused NDN/AI libraries.
\end{itemize}
\bluebox{This proposal defines those obligations. It does not label the remaining collaboration evidence as finished.}
\end{frame}

\begin{frame}{Discussion: Is This Scope Sufficient and Feasible?}
\begin{enumerate}
\item Are the three research questions specific enough to support defensible answers?
\item Are the authorization alternatives complete and fairly configured?
\item Are one UAV graph and one supported DI partition sufficient to test the abstraction?
\item Which optional features should be excluded if they delay the core evidence?
\end{enumerate}
\vspace{4mm}\bluebox{NDNSF: runtime service discovery and authorization, connected to explicit named-data execution rules.}
\end{frame}

\appendix
\begin{frame}{Backup: Discovery Configuration Is Not Framework Capability}
\small
\begin{tabularx}{\linewidth}{Xr}
\toprule Provider-transition configuration & Completed\\\midrule
NDNSF, no pre-registered Provider addresses & 357 / 357\\\midrule
Static-list NSC/gRPC configurations & 0 / 357\\\midrule
Pre-registered baseline variants & 357 / 357\\\bottomrule
\end{tabularx}
\vspace{5mm}\bluebox{This isolates the configured discovery boundary. Resolver/discovery integration can update candidates; it is not proof that gRPC or NSC cannot do so. [27,28]}
\end{frame}

\begin{frame}{Backup: Historical Loss Evidence Is Not a New Matched Test}
\small
\begin{tabularx}{\linewidth}{Xr}
\toprule Historical per-link loss & NDNSF completed requests\\\midrule
1\% & 599 / 600\\\midrule
3\% & 598 / 600\\\midrule
10\% (best recorded run) & 539 / 600\\\bottomrule
\end{tabularx}
\vspace{4mm}\bluebox{A selected run is not an aggregate estimate. Do not combine these numbers with another batch and label the result ``matched.'' [28]}
\vspace{2mm}\scriptsize $1-(1-\ell)^4$ assumes four independent equal-loss packet traversals; it is not service failure probability after retries.
\end{frame}

\begin{frame}{References: NDN and Authorization}
\scriptsize
\refentry{1}{L. Zhang et al., ``Named Data Networking,'' ACM SIGCOMM CCR, 2014.}
\refentry{2}{T. Li et al., ``A Brief Introduction to NDN Dataset Synchronization (NDN Sync),'' MILCOM, 2018.}
\refentry{3}{P. Moll et al., ``A Brief Introduction to State Vector Sync,'' NDN-0073, rev. 2, 2021.}
\refentry{4}{Z. Zhang et al., ``NAC: Automating Access Control via Named Data,'' MILCOM, 2018.}
\refentry{12}{K. Nichols, ``Lessons Learned Building a Secure Network Measurement Framework using Basic NDN,'' ACM ICN, 2019.}
\refentry{16}{Y. Yu et al., ``Schematizing Trust in Named Data Networking,'' ACM ICN, 2015.}
\refentry{19}{S. Dulal et al., ``Enhancing NAC-ABE to Support Access Control for mHealth Applications and Beyond,'' arXiv:2311.07299, 2023.}
\refentry{26}{NDN Project, ``Data Packet,'' NDN Packet Format Specification. docs.named-data.net/NDN-packet-spec/current/data.html (accessed Sept. 2026).}
\end{frame}

\begin{frame}{References: Invocation and Composition}
\scriptsize
\refentry{6}{C. Tschudin and M. Sifalakis, ``Named Functions and Cached Computations,'' IEEE CCNC, 2014.}
\refentry{7}{M. Krol and I. Psaras, ``NFaaS: Named Function as a Service,'' ACM ICN, 2017.}
\refentry{9}{M. Krol et al., ``RICE: Remote Method Invocation in ICN,'' ACM ICN, 2018.}
\refentry{10}{D. Meirovitch and L. Zhang, ``NSC---Named Service Calls, or a Remote Procedure Call for NDN,'' NDN-0074, rev. 1, 2021.}
\refentry{15}{CaSCON, SMARTCOMP 2025. Primary publication record: RIT, ``Barrios presents at IEEE SmartComp.'' Full feature-level comparison remains a research obligation.}
\refentry{17}{gRPC Authors, ``Core Concepts, Architecture and Lifecycle,'' gRPC documentation.}
\refentry{27}{gRPC Authors, ``Custom Name Resolution,'' grpc.io/docs/guides/custom-name-resolution/ (accessed Sept. 2026).}
\end{frame}

\begin{frame}{References: Tools, Evidence and Degree Planning}
\scriptsize
\refentry{22}{Mini-NDN, project documentation: minindn.ndn.today.}
\refentry{23}{PX4, ``Software-in-the-Loop Simulation,'' PX4 User Guide.}
\refentry{24}{NDN Project, ``Repo Protocol Specification,'' repo-ng documentation.}
\refentry{25}{ONNX, ``Open Neural Network Exchange Intermediate Representation Specification.''}
\refentry{28}{NDNSF, research-revision-audit.md, companion claim/evidence ledger, September 9, 2026. Contains campaign paths and unresolved acceptance gates.}
\refentry{29}{University of Memphis CS, ``Information for Computer Science Ph.D. Students,'' forms\_grad\_phd\_programinfo.pdf (accessed September 9, 2026).}
\refentry{30}{University of Memphis Graduate School, ``Graduation Information,'' Spring 2027 calendar (accessed September 9, 2026).}
\end{frame}
